%% ============================================================
%%  Section: Main Results  (2026-04-07 — addresses all [@@@ ...] comments)
%%  Depositor Discipline in Modern Banking
%%
%%  Changes from empirical_results_20260404_v2.tex:
%%    - Removed all [@@@ ...] placeholders and addressed each one.
%%    - Removed all \paragraph{} headers; content flows as prose.
%%    - Triple-difference folded into time deposits subsection (not standalone).
%%    - Opening paragraph shortened; "Why time deposits" condensed to one paragraph.
%%    - Controls paragraph removed; defers to Section~\ref{sec:identification}.
%%    - Economic magnitudes expressed as percent of sample mean, not dollar amounts.
%%    - Overconfident phrasing ("immune by construction", "cleanly validates",
%%      "central threat", "cleanest test") softened throughout.
%%    - Insured-deposit interpretation condensed; mechanical reclassification noted.
%%    - Removed the Post x Uninsured explanatory paragraph from triple-diff.
%%    - Financial consequences and sophistication sections shortened.
%%    - Summary subsection deleted.
%%    - NIM binary coefficient reported as 0.024 pp/quarter (consistent with table).
%%    - \label{sec:results:heterogeneity} added as cross-reference alias.
%%
%%  Cross-references created:
%%    sec:results
%%    sec:results:time_deposits
%%    sec:results:financial
%%    sec:results:sophistication
%%    sec:results:heterogeneity  (alias for sophistication)
%%
%%  Cross-references consumed:
%%    tab:small_bank_did_time_deposits
%%    tab:small_bank_triple_diff_time_deposits
%%    tab:small_bank_did_financial_outcomes
%%    fig:es_time_deps_cecl
%%    fig:es_financial_outcomes_cecl
%%    fig:es_sophistication_split
%%    sec:identification
%% ============================================================

\section{Main Results}
\label{sec:results}

We test whether the CECL Day-One disclosure triggered a disciplinary response from
uninsured depositors.  Three predictions follow from the theory: banks whose CECL
adjustment revealed larger latent credit losses should experience relative declines
in uninsured deposit balances after adoption, should face higher funding costs to
replace those balances, and should suffer lower profitability as a consequence.
The results confirm all three.

\subsection{Time Deposits}
\label{sec:results:time_deposits}

We focus on time deposits because they are the deposit category most susceptible
to disciplinary withdrawal.  Certificates of deposit and other fixed-term
instruments are held for investment rather than transaction purposes, so discipline
operates at the rollover margin: a depositor who declines to renew a maturing CD
is acting on a risk judgment without bearing the switching costs that complicate
responses in demand accounts \citep{calomiris1991role, iyer2016tale}.  Uninsured
deposits above the \$250{,}000 FDIC threshold are disproportionately held by
businesses, municipalities, and financially sophisticated households with both
the incentive and the capacity to act on new risk disclosures \citep{flannery1998using,
gee2022decision}.  Prior evidence for depositor discipline is strongest in
time-deposit data \citep{maechler2006dynamic, egan2017deposit}.

Table~\ref{tab:small_bank_did_time_deposits} presents difference-in-differences
estimates for uninsured and insured time deposits, using the specification in
Section~\ref{sec:identification}.  The interaction of the post-adoption indicator
with the binary High-CECL treatment (column~2) yields a coefficient of $-0.394$
percentage points of assets (standard error $0.182$, $p < 0.05$), roughly 7
percent of the sample mean for uninsured time deposits.  The continuous
specification (column~1) is consistent: each additional percentage point of CECL
exposure reduces uninsured time deposits by $0.072$ percentage points of assets
($p < 0.05$).

Insured time deposits respond in the opposite direction.  The High-CECL
coefficient in column~4 is $+0.712$ percentage points ($p < 0.05$), indicating
that high-exposure banks attracted more insured balances after adoption.  The
continuous estimate (column~3, $+0.070$ pp) is positive but imprecise.  The
insured response is broadly consistent with banks substituting toward
insurance-protected funding after losing uninsured deposits, as theory predicts
\citep{calomiris1991role}.  The positive insured coefficient may also partly
reflect mechanical reclassification as some depositors reduce balances below the
\$250{,}000 threshold rather than exit the bank entirely.

Figure~\ref{fig:es_time_deps_cecl} plots event-study coefficients for a
ten-quarter window on either side of CECL adoption ($k = -1$ omitted).
Pre-adoption coefficients for uninsured time deposits are clustered near zero
across all ten lead quarters and are jointly indistinguishable from the reference
period, consistent with the parallel-trends assumption.  The post-adoption break
emerges within the first adoption quarter and persists through the end of the
sample horizon, while the insured series moves in the opposite direction with
comparable timing.  The absence of pre-adoption divergence in either series
argues against the interpretation that the post-adoption pattern reflects a
differential trend predating the disclosure.

A remaining concern is that high-CECL banks may have faced bank-specific funding
pressures after 2023 that depressed their entire time-deposit portfolio
independently of depositor discipline.  If CECL merely compressed aggregate
deposit demand at high-exposure institutions, insured and uninsured balances
would co-move, and the difference-in-differences estimate would overstate
discipline effects \citep{covitz2004market}.  To address this, we stack the
insured and uninsured time-deposit series in a single panel and estimate a
triple difference:
\begin{equation*}
\begin{split}
  y_{idt}
  = \alpha_{it}
  + \beta_1 \cdot \mathrm{Uninsured}_d
  + \beta_2 \cdot \mathrm{Post}_{t} \times \mathrm{Uninsured}_d
  + \beta_3 \cdot \mathrm{CECL}_i \times \mathrm{Uninsured}_d \\
  \quad + \beta_4 \cdot \mathrm{CECL}_i \times \mathrm{Post}_{t} \times \mathrm{Uninsured}_d
  + \varepsilon_{idt},
\end{split}
\end{equation*}
where $d \in \{\text{uninsured}, \text{insured}\}$ indexes deposit type and
$\alpha_{it}$ are bank $\times$ quarter fixed effects.  These absorb the full
vector of bank-specific time-varying shocks in each period, so identification
comes from the differential movement of uninsured relative to insured deposits
within the same bank and the same quarter.  The triple-interaction coefficient
$\hat{\beta}_4$ captures the within-bank composition shift.

Table~\ref{tab:small_bank_triple_diff_time_deposits} reports the estimates.  The
triple interaction in the binary specification (column~2) is $-1.252$ (standard
error $0.340$, $p < 0.01$): relative to their own pre-adoption composition and
relative to low-CECL banks, high-CECL banks experienced a 125-basis-point
widening of the insured-to-uninsured gap as a share of assets after adoption.
The continuous analog (column~1, $-0.177$, SE $0.058$, $p < 0.01$) confirms a
smooth dose-response relationship across the CECL distribution rather than a
threshold phenomenon concentrated in the extreme tail.  Because bank $\times$
quarter fixed effects absorb any level shift in total time-deposit demand at each
bank in each quarter, these estimates are not consistent with any story in which
CECL simply depressed aggregate funding at high-exposure institutions.

\subsection{Financial Consequences}
\label{sec:results:financial}

Table~\ref{tab:small_bank_did_financial_outcomes} reports difference-in-differences
estimates for interest expense, return on assets, and net interest margin.
High-CECL banks paid $0.021$ percentage points more per quarter in interest
expense per dollar of assets after adoption (column~2, $p < 0.01$), an
annualized increase of approximately 8--9 basis points, or about 9 percent of
the sample quarterly mean.  The continuous specification implies that each
additional percentage point of CECL exposure raises quarterly interest expense
by 0.37 basis points of assets (column~1, $p < 0.05$).  Having lost lower-cost
uninsured balances, high-CECL banks replaced them with insured time deposits at
rates sufficient to attract depositors who would not otherwise bear uninsured
credit risk at these institutions, the price channel of depositor discipline
\citep{nier2006market, billett1998cost}.

The higher funding cost was not offset elsewhere.  Return on assets fell by
$0.038$ percentage points per quarter for high-CECL banks (column~4, $p < 0.01$),
an annualized drag of about 15 basis points and roughly 13 percent of the sample
mean.  Net interest margin contracted by $0.024$ percentage points per quarter
(column~6, $p < 0.05$).  The joint decline in ROA and NIM is consistent with
higher liability costs not being passed through to borrowers, in line with
community banks typically facing competitive local loan markets where
short-run asset-side repricing is limited \citep{hannan1988bank}.
Continuous estimates scale proportionally: each percentage point of CECL equity
shock reduces quarterly ROA by $0.009$ percentage points ($p < 0.01$) and
quarterly NIM by $0.005$ percentage points ($p < 0.05$).

Figure~\ref{fig:es_financial_outcomes_cecl} plots event-study coefficients for
all three financial outcomes.  Pre-adoption leads are flat across both treatment
specifications, and the post-adoption divergence for ROA and NIM persists through
the end of the sample horizon.

\subsection{Heterogeneity by Depositor Sophistication}
\label{sec:results:sophistication}
\label{sec:results:heterogeneity}

The disciplinary mechanism requires depositors who are willing and able to act
on new risk information.  We test whether the response is amplified in markets
with more financially sophisticated depositor bases, splitting the sample at the
75th percentile of a deposit-weighted branch-market sophistication index
constructed from 2022 FDIC Summary of Deposits matched to ZIP-code-level
American Community Survey data \citep{ACS2022}.  Banks at or above this threshold
are classified as operating in high-sophistication markets.

Figure~\ref{fig:es_sophistication_split} presents event-study estimates separately
for high- and low-sophistication banks.  Panel~A shows that the uninsured deposit
response is concentrated in high-sophistication markets: the event-study
coefficient for that group drifts steadily downward after adoption, reaching
approximately $-1.5$ to $-2$ percentage points of assets by quarters $+9$ to $+11$,
while the low-sophistication group remains near zero throughout.  The gradual
buildup explains why the pooled triple-interaction coefficient ($\mathrm{Post}
\times \mathrm{High\,CECL} \times \mathrm{High\,Sophistication} = -0.022$,
SE $= 0.468$) is statistically insignificant: the post-period average pools early
quarters before the separation is visible with later quarters where the gap is
large.  The pattern is consistent with sophisticated depositors redirecting
balances gradually as time deposits mature, rather than breaking term contracts
early \citep{iyer2016tale}.

The sharpest statistical signal from the sophistication split appears in the price
channel.  Panel~C shows that high-CECL banks in high-sophistication markets paid
a distinctly steeper increase in interest expense after adoption; the
triple-interaction coefficient is $+0.044$ percentage points per quarter
($p = 0.020$), about 17--18 additional basis points annually on top of the
pooled effect.  Panel~D confirms that the additional funding cost translated into
larger profitability losses: the triple-interaction for ROA is $-0.076$ percentage
points per quarter ($p = 0.010$), an annualized drag of approximately 30
additional basis points.  Panel~E shows a directionally consistent but
statistically imprecise effect on NIM ($-0.031$, $p = 0.175$).  Panel~B shows
that high-sophistication banks also attracted fewer insured deposits than their
low-sophistication counterparts after adoption ($-0.965$, SE $= 0.677$,
$p = 0.154$), consistent with sophisticated depositors being more reluctant to
serve as replacement insured funding at institutions whose credit quality has
deteriorated.  Together, these results suggest a two-channel discipline
mechanism: sophisticated depositors extract higher rates immediately upon
observing the CECL disclosure, then gradually reduce balances as term instruments
mature \citep{egan2017deposit}.
